\section{Problem formulation and method}\label{sec:method}

\subsection{Objective}
Given observations $(x_i,y_i)$, $i=1,\dots,n$, with $x_i\in\R^d$ and
$y_i\in\{-1,+1\}$, we minimize
\begin{equation}\label{eq:obj}
  F(w)=\frac1n\sum_{i=1}^n \log\bigl(1+e^{-z_i}\bigr)+\frac\lambda2\norm{w}_2^2,
  \qquad z_i=y_i\,x_i^\top w,
\end{equation}
with $\lambda=10^{-6}$ unless stated otherwise and no separate intercept. Writing
$X\in\R^{n\times d}$ for the design matrix and $\sigma(t)=1/(1+e^{-t})$,
\begin{equation}\label{eq:derivs}
  \nabla F(w)=-\frac1n X^\top\bigl(y\odot\sigma(-z)\bigr)+\lambda w,\qquad
  \nabla^2F(w)=\frac1n X^\top D X+\lambda I,\quad D=\operatorname{diag}\bigl(\sigma(z_i)(1-\sigma(z_i))\bigr).
\end{equation}
Since $\lambda>0$, $F$ is $\lambda$-strongly convex and has a unique minimizer $w^*$.
The weights $D_{ii}\le 1/4$ vanish for large $|z_i|$, so the Hessian inherits and
can amplify the ill-conditioning of $X^\top X/n$: its smallest eigenvalue can be as
small as $\lambda$.

\subsection{Stable evaluation}
The loss term is evaluated as $\log(1+e^{-z})=\texttt{logaddexp}(0,-z)$, which
computes $\max(0,-z)+\log(1+e^{-|z|})$ and never exponentiates a positive number.
The sigmoid uses $\sigma(t)=1/(1+e^{-t})$ for $t\ge0$ and $e^{t}/(1+e^{t})$ for
$t<0$. Both forms are exact reformulations, so they change nothing when the naive
formulas do not overflow. For the ablation in \cref{sec:results} we also implement the
naive formulas $\log(1+e^{-z})$ and $1/(1+e^{-z})$.

\subsection{Search directions and line search}
Let $g_k=\nabla F(w_k)$ and $H_k=\nabla^2F(w_k)$. The methods compared use
\begin{align*}
  \text{GD:}&\quad p_k=-g_k, &
  \text{Newton:}&\quad H_kp_k=-g_k \ \text{(dense solve)},\\
  \text{Newton--CG:}&\quad H_kp_k\approx-g_k \ \text{(CG)}, &
  \text{ASH-LR (phase 2):}&\quad (H_k+\mu_kI)p_k\approx-g_k \ \text{(CG)}.
\end{align*}
L-BFGS uses the two-loop recursion with memory $m=10$ and skips curvature pairs
with $s^\top y\le10^{-10}\norm{s}\norm{y}$. Except for GD with fixed step $1/L$,
$L=\norm{X}_2^2/(4n)+\lambda$, all line searches use Armijo
backtracking~\cite{armijo1966minimization}: starting from $\alpha=1$, the step is
halved until $F(w_k+\alpha p_k)\le F(w_k)+c\,\alpha\, g_k^\top p_k$ with $c=10^{-4}$,
for at most 40 trials.

\subsection{Inexact solves with Hessian-vector products}
The product $(H_k+\mu I)v=\frac1n X^\top\bigl(D_k(Xv)\bigr)+(\lambda+\mu)v$
costs two passes over the data and never forms $H_k$. CG is run on
$(H_k+\mu_kI)p=-g_k$ from $p=0$ and stopped when the residual
$r=-g_k-(H_k+\mu_kI)p$ satisfies
\begin{equation}\label{eq:forcing}
  \norm{r}\le\eta_k\norm{g_k},\qquad \eta_k=\min\bigl(0.5,\sqrt{\norm{g_k}}\bigr),
\end{equation}
or after $d$ iterations. Because $H_k+\mu_kI$ is positive definite, every CG
iterate is a descent direction~\cite{nocedal2006numerical}.

\subsection{Switching rule}
ASH-LR starts with Armijo gradient steps and switches, once and for all, to the
damped Newton phase as soon as either
(i) $\norm{g_k}\le 0.1\,\norm{g_0}$, or
(ii) at least five gradient steps have been taken and the last step achieved a relative
decrease $(F(w_{k-1})-F(w_k))/\max(1,|F(w_k)|)\le10^{-3}$.
The idea is that cheap gradient steps make fast progress while far from the solution
and while the curvature model is unreliable, and that Newton steps take over when
gradient steps stall on the low-curvature directions.

\subsection{Damping by a ratio test}
For a trial direction $p$ the quadratic model predicts the decrease
\begin{equation}\label{eq:pred}
  \mathrm{pred}(p)=-\bigl(g_k^\top p+\tfrac12 p^\top H_kp\bigr).
\end{equation}
With an inexact solve, $(H_k+\mu_kI)p=-g_k-r$, so
$p^\top H_kp=-p^\top g_k-p^\top r-\mu_k\norm{p}^2$ is available from CG quantities
without an extra Hessian-vector product. The ratio
$\rho=\bigl(F(w_k)-F(w_k+p)\bigr)/\mathrm{pred}(p)$ measures how well the model
predicted the actual decrease. The step is accepted with $\alpha=1$ if
$\rho\ge10^{-4}$; otherwise it is rejected and the direction is recomputed with
larger damping. The damping parameter follows the gradient-regularized rule
\begin{equation}\label{eq:mu}
  \mu_k=\operatorname{clip}\bigl(\sqrt{\theta_k\norm{g_k}},\,10^{-12},\,10^{8}\bigr),
\end{equation}
and the ratio test adapts $\theta_k$: $\theta\leftarrow4\theta$ if $\rho<0.25$,
$\theta\leftarrow\theta/2$ if $\rho>0.75$, unchanged otherwise. At the switch
$\theta$ is initialized so that $\mu=10^{-3}$. Because $\mu_k\to0$ as
$\norm{g_k}\to0$, the method becomes an undamped inexact Newton method near the
solution, whereas in the Milestone~2 rule $\mu$ was only halved after each full
step and could remain large for many iterations. As a safeguard, after 30 consecutive
rejections the last direction is used with Armijo backtracking.

\Cref{alg:ash} summarizes the method. Every accepted step decreases $F$: gradient
and fallback steps satisfy the Armijo condition, and ratio-test steps satisfy
$F(w_k)-F(w_{k+1})\ge10^{-4}\,\mathrm{pred}(p)>0$.

\begin{algorithm}[t]
\caption{ASH-LR}\label{alg:ash}
\begin{algorithmic}[1]
\Require $w_0$, tolerance $\varepsilon$, $\mu_0=10^{-3}$
\State phase $\gets$ GD; $k\gets0$
\While{$\norm{g_k}>\varepsilon$ and $k<k_{\max}$}
  \If{phase $=$ GD}
    \State $p\gets-g_k$; $\alpha\gets$ Armijo$(w_k,p)$; $w_{k+1}\gets w_k+\alpha p$
    \If{$\norm{g_{k+1}}\le0.1\norm{g_0}$ \textbf{or} (GD steps $\ge5$ \textbf{and} rel.\ decrease $\le10^{-3}$)}
      \State phase $\gets$ Newton; $\theta\gets\mu_0^2/\norm{g_{k+1}}$
    \EndIf
  \Else
    \For{$j=1,\dots,30$}
      \State $\mu\gets\operatorname{clip}(\sqrt{\theta\norm{g_k}})$; solve $(H_k+\mu I)p=-g_k$ by CG to tolerance \eqref{eq:forcing}
      \State $\rho\gets(F(w_k)-F(w_k+p))/\mathrm{pred}(p)$
      \State $\theta\gets4\theta$ if $\rho<0.25$; \quad $\theta\gets\theta/2$ if $\rho>0.75$
      \If{$\rho\ge10^{-4}$} \State $w_{k+1}\gets w_k+p$; \textbf{break} \EndIf
    \EndFor
    \State \textbf{if} no step accepted: $w_{k+1}\gets w_k+\text{Armijo}(w_k,p)\,p$
  \EndIf
  \State $k\gets k+1$
\EndWhile
\end{algorithmic}
\end{algorithm}

\subsection{Cost model}
We measure work in \emph{data passes}, one pass being a product with $X$ or
$X^\top$ ($\approx nd$ multiply-adds). A function value costs one pass, a gradient
at a point whose margins are cached one more pass, a Hessian-vector product two
passes, forming $X^\top DX$ costs $d$ passes, and a dense solve $d^2/(3n)$ passes.
Passes are machine-independent and are the primary efficiency measure; wall time is
reported as well, because dense matrix--matrix products run much closer to peak
speed than the matrix--vector products used by CG, so equal passes do not mean
equal time.
